\chapter{Vendetta in SOL maggiore}




Laura, seduta sulla poltroncina \textit{for friends}, aiuta l’amica a sbrigare un po’ della corrispondenza mentre lei è impegnata in altre faccende che le stanno più a cuore.

\par\medskip
\begin{center}
  \includegraphics[width=.75\linewidth]{image1.png}
\end{center}

Una lettera è un messaggio scritto di pugno e consegnato nella buchetta di casa per mezzo di un servizio postale.
Per molte persone ha un significato simbolico: indica l’appartenenza al gruppo di chi non ha dimenticato l’origine terrestre degli esseri umani.
In molti si sono convinti che la carta non sia il male del mondo ed esprimono con le parole il loro dolore per il pianeta che soffre, scrivendole con la matita, sulla carta.
Una matita fatta di legno sulla carta di cellulosa.
Parole fatte di albero, spedite a destinatari incaricati di sorreggere il peso dell’angoscia collettiva.

\begin{verse}\itshape
«Cara Caterina,\\
ti scrivo con le mani sporche di terra e il cuore in fiamme.\\
Ogni giorno mi sveglio con l’ansia che il cielo sia un po’ più basso,\\
che il vento porti con sé un alto grido soffocato di una specie che non c’è più.\\
Eppure continuo a lottare, perché se ci sei tu,\\
tu che hai il coraggio che io non ho…»
\end{verse}

«Vuoi che continui, Cate? Non mi sembra aggiunga nulla di nuovo alla discussione…»\\
«Ma sì dai, mi sembra così angosciata, poverina…»
\par\medskip
Caterina è a pochi decimetri da Laura, davanti all’armadio
\par\medskip
\begin{center}
  \includegraphics[width=.75\linewidth]{image3.png}
\end{center}

che sceglie i vestiti da mettere in valigia.

Laura riprende a leggere ma il settanta per cento delle parole sfugge alla sua attenzione focalizzata sui movimenti mimetici di Alice che seduta sul letto sul letto di fronte a lei armeggia con un vecchio portatile ostentando disinvoltura.\\

Caterina è su di giri, ma esausta. Tra biglietti introvabili, alberghi al completo e impegni incastrati all’ultimo, è stato un viaggio difficile da organizzare.

Mancano trentasei ore al volo che la condurrà negli Stati Uniti ed è quasi tutto a posto, ma  non sa ancora che sua sorella le sta cucinando una bella vendetta da gustare calda anziché fredda.

Caterina  va a New York lasciando Alice a casa con suo padre e sua madre: \textit{proprio ora che potremmo passare una settimana insieme} rimugina Alice continuando a scrivere sulla tasiera.  
Certo, Alice comprende che il lavoro è importante e se si trattasse  solo di lavoro, non le brucerebbe così tanto.

Ma il problema è un altro: \textit{lei ci va per Mark}.  
Figurati, se non ci fosse stati lui, avrebbe sicuramente preferito passare l’estate con lei.  
In fondo, di corsi di aggiornamento ce ne sono tanti nel mondo.  
\textit{Lei ci va per Mark.}

\par\medskip
Così, tranquilla, con il PC in mano e le gambe incrociate sul letto, mostra uno sguardo innocente, perfettamente a suo agio nell’ambiente. Intorno a lei, però, non c’è la sua stanza.  
È la camera di Caterina: ordinata, spaziosa, un piccolo ecosistema.

Dalla camera di una ragazza si può capire molto di lei.

\par\medskip
\begin{center}
  \includegraphics[width=.75\linewidth]{image2.png}
\end{center}

L'ampio spazio al centro racconta del suo cuore e della libertà di cui ha bisogno. I mobili accostati alle pareti testimoniano, involontari, il suo bisogno di organizzazione e sicurezza, fisica e affettiva, quella che Alice continua a sfidare creando una tensione forte e percepibile, ma senza spiegarne la ragione.
Laura però vuole capire,  si avvicina per sbirciare cosa stia facendo, ma Alice, con una piccola rotazione, si sottrae al suo sguardo.

«Scusami, sai!» la rimprovera dell'indiscrezione.  
Poi incolla il testo nella sezione \textit{Lyrics}:
\pagebreak
\par\medskip
\textbf{Lyrics}
\begin{quote}\itshape
She leaves and I stay\\
like a folder left open\\
half full, half erased\\
I blink, and she’s already gone
\end{quote}

«Qui capirai quanto ci sono rimasta male…»

\par\medskip
\textbf{Style Description}
\begin{quote}\itshape
Bilingual emotional synth-pop with ambient textures, cinematic flow, and AI female vocals.\\
Slow build. Glitchy, dreamlike, bittersweet. Like diary pages sung in code.
\end{quote}

«Lasciala perdere, è arrabbiata con me.»

«Non sono arrabbiata, solo che voglio finire una cosa.»

«Ha ragione, sono io che sono troppo curiosa.  
Però si è fatto un po’ tardi, è ora che vado a prendere Valentina.»

«No, aspetta solo un attimo, ti prego.»

«Ora la scarico, colleghiamo le casse e voilà!»

\par\medskip
Le note di pianoforte sintetico arpeggiano velocemente una nenia in Sol maggiore.  
Poi, con un respiro sussurrato, inizia il primo verso:

\par\medskip
\textbf{Verse 1}
\begin{quote}\itshape
She leaves and I stay\\
like a folder left open\\
half full, half erased\\
I blink, and she’s already gone
\end{quote}

Questa è la scena prima della crisi…

\par\medskip
\begin{center}
  \includegraphics[width=.75\linewidth]{ccsf.png}
\end{center}

\textbf{Verse 2}
\begin{quote}\itshape
She goes to New York\\
I stay with a lamp shaped like a heart\\
plastic love, three settings\\
warm, cold, ambient — mine is blinking
\end{quote}

«E qui voglio vederti piangere!»

\begin{quote}\itshape
I'm not angry, I'm just here
\end{quote}
Caterina continua ad ordinare le cose da mettere in valigia, mentre una lacrima le riga il viso. Laura la vede scendere piano fino alla guancia, lì Caterina la ferma, pulendosi con i polpastrelli. «Certo che le semplifichi proprio la partenza, tu.» Le parole di Laura sciolgono Caterina che appoggia l’asciugacapelli e si avvicina alla sorella per abbracciarla.

«Non ci provare!» le urla Alice.

Caterina non reagisce. È abituata.  
Alice si libera le ginocchia dal PC e lo poggia sul letto.

«Io esco, mangio qualcosa con le mie amiche.»

Caterina si asciuga gli occhi. Ormai è inutile simulare.

«Va bene, ma torna a casa per le dieci. È l’ultima sera che passiamo insieme.»\\
Laura si alza e si sistema i vestiti:
«Devo andare anch’io, Cate.»
\\
«Va bene Laura, grazie di essere venuta.»
\\
«Ciao Cate. Ciao Alice.»
\\
Prima di chiudere la porta, Laura ha un attimo di esitazione.  
C’era ancora una cosa, anzi, era il motivo principale per cui aveva raggiunto Caterina.

«Ma non sai ancora nulla del visto?»\\

«Non preoccuparti, Laura. Vedrai che non ci saranno problemi.  
In ogni caso, se ci fossero difficoltà ti chiamo.»

Saluta l’amica con un bacio e corre a prendere Vale.  
Per fortuna manca ancora qualche anno alla sua adolescenza.
